\documentclass{article}
\usepackage[UTF8]{ctex}
\usepackage{amsmath,amssymb,amsfonts}
\usepackage{enumitem}
\usepackage{geometry}
\geometry{a4paper, margin=1in}
% === 错题本排版优化（不参与题目标识） ===
\setlength{\parindent}{0pt}
\setlength{\parskip}{0.35em plus 0.10em minus 0.08em}
\setlength{\emergencystretch}{3em}
\setlist{itemsep=0.25em, topsep=0.40em, parsep=0pt}
\allowdisplaybreaks[2]
% === 错题本排版优化结束 ===
\begin{document}
\section*{西南交通大学2025年高等代数真题}
\section*{162．西南交通大学2025年高等代数试卷}

一．填空题．

\subsection{题目1}
1．设二次型 \(f\left(x_{1}, x_{2}, x_{3}, x_{4}\right)\) 的正惯性指数为 1 ，且 \(f\left(x_{1}, x_{2}, x_{3}, x_{4}\right)\) 的矩阵 \(A\) 满足 \(A^{2}-2 A=3 E\) ，则 \(f\left(x_{1}, x_{2}, x_{3}, x_{4}\right)\) 的符号差为 \(\_\_\_\_\)。

\subsection{题目2}
2．设矩阵 \(A=\left(\begin{array}{lll}1 & 2 & 1 \\ 3 & 4 & a \\ 1 & 2 & 2\end{array}\right)\) ，矩阵 \(B\) 满足 \(A B=A-B+E\) ，其中 \(a\) 为常数，\(E\) 为单位矩阵，且 \(B \neq E\)，则 \(a=\) \(\_\_\_\_\)．

\subsection{题目3}
3．设 \(A, B\) 为2025阶正交矩阵，齐次线性方程组 \((A-B) X=0\) 的解空间维数为4，则 \(A B^{-1}\) 的行列式 \(\operatorname{det}\left(A B^{-1}\right)=\) \(\_\_\_\_\)。

\subsection{题目4}
4．设 \(f(x)=x^{5}+2 x^{4}-2 x^{3}-8 x^{2}-7 x-2\) ，则 \(f(x)\) 在复数域上不可约分解式为 \(f(x)=\) \(\_\_\_\_\)．

\subsection{题目5}
5．将 \(n\) 阶全体实对称矩阵合同关系进行分类（ \(A\) 与 \(B\) 合同当且仅当存在可逆的实矩阵 \(P\) 使得\(P^{\mathrm{T}} A P=B\) ，其中 \(P^{\mathrm{T}}\) 为 \(P\) 的转置），则一共有 \(\_\_\_\_\)类．

\section*{二．解答题。}

\subsection{题目6}
1．设
\[
\alpha_{1}=\left(\begin{array}{c}
1 \\
2 \\
-1 \\
0 \\
4
\end{array}\right), \alpha_{2}=\left(\begin{array}{c}
-1 \\
3 \\
2 \\
4 \\
1
\end{array}\right), \alpha_{3}=\left(\begin{array}{c}
2 \\
9 \\
-1 \\
4 \\
13
\end{array}\right)
\]
\(W=L\left(\alpha_{1}, \alpha_{2}, \alpha_{3}\right)\) 是 \(\alpha_{1}, \alpha_{2}, \alpha_{3}\) 生成的线性空间 \(\mathbb{R}^{6}\) 的子空间，其中 \(\mathbb{R}\) 为实数域．
（1）求以 \(W\) 为解空间的齐次线性方程组。
（2）求以 \(\{\eta+\alpha \mid \alpha \in W\}\) 为解集的非齐次线性方程组，其中 \(\eta=\left(\begin{array}{c}1 \\ 2 \\ 1 \\ 2 \\ 1\end{array}\right)\) ．

\subsection{题目7}
2．设 \(A\) 是一个 \(n\) 阶复方阵且秩 \((A)=1, E\) 为 \(n\) 阶单位矩阵。求 \(A-E\) 的 Jordan 标准形。

\subsection{题目8}
3．设 \(V\) 是实数域上所有 \(n\) 阶实对称矩阵构成的线性空间，对任意的 \(A, B \in V\) ，定义 \((A, B)=\operatorname{Tr}(A B)\) ， \(\operatorname{Tr}(A B)\) 表示乘积矩阵 \(A B\) 的迹。

（1） 证明：\(V\) 关于 \((A, B)\) 形成一个欧氏空间。
（2） 令 \(S=\{A \in V \mid \operatorname{Tr}(A)=0\}\) ，求 \(S\) 及其正交补 \(S^{\perp}\) 的一组基与维数。

\subsection{题目9}
4．设 \(n\) 阶方阵 \(A\) 满足 \(A^{3 m}+A+E=O\) ，其中 \(E\) 为 \(n\) 阶单位矩阵，\(m\) 为正整数。证明：\(A^{2}+A+E\)是可逆矩阵，并求其逆矩阵。

\subsection{题目10}
5．设 \(p_{1}, p_{2}, \cdots, p_{n}\) 是 \(n\) 个互不相同的素数，证明：\(\sqrt[n]{p_{1} p_{2} \cdots p_{n}}\) 必是无理数，其中 \(n\) 是正整数，且\(n>1\) ．

\subsection{题目11}
6．设 \(A, B\) 为 \(n\) 阶方阵，齐次线性方程组 \(A X=0\) 和 \(B X=0\) 同解且它们的基础解系中包含 \(m\) 个线性无关的解向量。证明：秩 \((B-A) \leq n-m\) 。

\subsection{题目12}
7．设 \(A\) 为 \(n\) 阶实方阵，且 \(E-A\) 的复特征值模长均小于 1 ，其中 \(E\) 表示 \(n\) 阶单位矩阵。证明：\(0<\operatorname{det}(A)<2^{n}\) ，这里 \(\operatorname{det}(A)\) 表示 \(A\) 的行列式。

\subsection{题目13}
8．设 \(A=\left(\begin{array}{ccc}1 & 0 & 1 \\ 4 & 1 & 3 \\ -3 & -1 & -2\end{array}\right), B=\left(\begin{array}{ccc}1 & 0 & 1 \\ 0 & 1 & -1 \\ 0 & 0 & 0\end{array}\right)\) ．问：

（1） 是否存在可逆矩阵 \(P\) ，使得 \(P A=B\) ？若存在，求出这样的一个矩阵 \(P\) ．
（2） 求满足 \(X A=B\) 的所有 3 阶方阵 \(X\) 。

\subsection{题目14}
9．设 \(V_{1}, V_{2}\) 是 \(n\) 维欧氏空间 \(V\) 的子空间，且 \(V_{1}\) 的维数小于 \(V_{2}\) 的维数．证明：在 \(V_{2}\) 中必存在一个非零向量正交于 \(V_{1}\) 中的一切向量。
\end{document}
